%<<CONT
% !PRINTED 228
\weberpage{228}
ist aber, wenn wir
\[
\Phi =
\begin{vmatrix}
1, & F_0(x_1), & F_1(x_1) & \ldots & F_{n-2}(x_1) \\
1, & F_0(x_2), & F_1(x_2) & \ldots & F_{n-2}(x_2) \\
\multicolumn{5}{c}{\cdots\cdots\cdots\cdots\cdots\cdots\cdots\cdots} \\
1, & F_0(x_n), & F_1(x_n) & \ldots & F_{n-2}(x_n)
\end{vmatrix}
\]
setzen, nach (11)
\weq{12}{\Phi^2 = n \varDelta.}
Beachtet man nun die Ausdrücke (9) in \S.\,68 für die Functionen
$F_0, F_1 \ldots F_{n-2}$, so ergiebt sich leicht nach dem Satze,
dass man in einer Determinante eine mit einem beliebigen Factor
multiplicirte Colonne zu einer anderen Colonne addiren kann
(\S.\,22, VII), für $\Phi$ der Ausdruck
\[
a_0^{n-1}
\begin{vmatrix}
1, & x_1, & x_1^2 & \ldots & x_1^{n-1} \\
1, & x_2, & x_2^2 & \ldots & x_2^{n-1} \\
\multicolumn{5}{c}{\cdots\cdots\cdots\cdots\cdots\cdots} \\
1, & x_n, & x_n^2 & \ldots & x_n^{n-1}
\end{vmatrix},
\]
dessen Quadrat nach \S.\,46 die Discriminante von $f(x)$ ist. Hiernach
haben wir also nach (12) den wichtigen Satz:

\gesperrt{Die Determinante der Bezoutiante von $f(x)$ ist der
$n^{\mathrm{te}}$ Theil der Discriminante von $f(x)$.}

Die Berechnung der Bezoutiante hat nach unseren Formeln
gar keine Schwierigkeit mehr und gestaltet sich ziemlich einfach.
Für $n = 3$ erhält man das schon aus den Formeln des
\S.\,71 zu schliessende Resultat
\weq{13}{B = - \frac{2}{3} H(t_1, t_0).}
Für die Berechnung der Bezoutiante der biquadratischen
Form wollen wir zur Veranschaulichung des Vorhergehenden
das Formelsystem vollständig aufführen, indem wir die Durchführung
der wenig längeren Rechnung für die Form fünften
Grades dem Leser überlassen.

Es ist für die Form vierten Grades nach (7) und \S.\,72, (9)
\[
f(x) = a_0 x^4 + a_1 x^3 + a_2 x^2 + a_3 x + a_4
\]
\[
\begin{array}{rcl|r}
E_{0,1} &=& -a_2 t_2 - a_3 t_1 - a_4 t_0 & 4 a_0 \tau_2 \\
E_{1,1} &=& a_1 t_2 & 3 a_1 \tau_2 \\
E_{2,1} &=& a_0 t_2 + a_1 t_1 & 2 a_2 \tau_2 \\
E_{3,1} &=& a_0 t_1 + a_1 t_0 & a_3 \tau_2 \\
\hline
\end{array}
\]
%CONT>>
